\chapter{Системийн зохиомж ба арга зүй}

\section{Системийн ерөнхий микросервис архитектур}
Энэхүү системийн ерөнхий архитектур нь микросервисийн зарчимд тулгуурлан үүсгэгдсэн. Энэхүү хандлага нь системийн бүрэлдэхүүн хэсгүүдийг хооронд нь сул хамааралтай (loosely coupled) байлгаж, тус бүрийг бие даан өргөтгөх, шинэчлэх, хянах боломжийг олгодог. Систем нь үндсэндээ Мэдээлэл цуглуулах давхарга (Гар утасны аппликейшн), Үйлчилгээний давхарга (FastAPI Backend API), ба Өгөгдлийн болон Оюун ухааны давхарга (MongoDB ба YOLOv8) гэсэн гурван үндсэн компонентоос бүрдэнэ.

\begin{enumerate}
    \item \textbf{Хэрэглэгчийн гар утасны аппликейшн (Flutter Frontend):} Хэрэглэгч (дэлгүүрийн менежер, түгээгч) нь уг гар утасны аппликейшнийг ашиглан дэлгүүрийн лангууны зургийг дарж эсвэл утасны санах ойгоос сонгож сервер рүү илгээнэ. Мөн серверийн буцааж өгсөн танилтын үр дүн, зөрүүтэй барааны тоо хэмжээ, аудитын түүх, статистик тайланг графикаар харна. Аппликейшн нь UI/UX -ийн орчин үеийн стандартыг хангасан байдлаар хийгдсэн бөгөөд сүлжээ тасарсан үед ч өөрийн локал санах ойд түр хадгалах боломжтойг системээс шаарддаг.
    \item \textbf{Backend сүлжээний үйлчилгээ (FastAPI):} Аппликейшнаас ирсэн HTTP multipart/form-data хүсэлтийг хүлээн авч, зургийг түр санах ойд хадгалан Машин сургалтын модуль руу дамжуулах төв зангилаа болж ажиллана. Мөн MongoDB өгөгдлийн сантай харьцаж шинэ бараа бүртгэх, устгах, засварлах CRUD (Create, Read, Update, Delete) үйлдлээр зохицуулна. FastAPI нь Python-ийн Pydantic санг хэрэглэн өгөгдлийн бүтцийн баталгаажуулалт (data validation) болон автомат OpenAPI (Swagger) баримтжуулалтыг хийдэг өндөр хүчин чадалтай фреймворк юм.
    \item \textbf{Компьютерын харааны модуль (ProductDetector):} Бэлтгэгдсэн YOLOv8 гүн сургалтын загварыг санах ойд байнгын ачаалан байрлуулсан (in-memory) үйлчилгээ байна. Зураг ирмэгц урьдчилан боловсруулалт хийгдэж (resize, normalization), загварт орж барааны ангилал, өөртөө итгэх хувь (confidence score) болон байршлыг цэгэн солбицлоор (bounding box) илрүүлэн гаргана. Ингээд үр дүнгээ JSON формат руу хөрвүүлж буцаана.
\end{enumerate}

\section{YOLOv8 гүн сургалтын загварын архитектур ба математик үндэслэл}
YOLOv8 загвар нь Ultralytics багаас гаргасан бөгөөд өмнөх YOLOv5-аас илүү сайжруулагдсан архитектуртай. Системд бид \texttt{yolov8n.pt} (Nano) буюу хамгийн бага параметр бүхий хувилбарыг сонгон ашиглах нь үүлэн сервер дээр эсвэл дотоод серверт тооцоолол хийхэд нөөцийг хэмнэх, хурд болон нарийвчлалын төгс тэнцвэртэй байдлыг хадгалах зохистой шийдвэр болно. YOLOv8 загварын дотоод бүтцийг авч үзвэл үндсэндээ:
\begin{itemize}
    \item \textbf{Backbone (Сээр нуруу):} CSPDarknet53 архитектурын сайжруулсан хувилбарыг ашигладаг. Оролтын зургаас үндсэн онцлог шинж чанаруудыг (өнгө, хэлбэр, зураас) татаад авдаг хэсэг. Шинэ загварт \textit{C3} (Cross Stage Partial) модулийг \textit{C2f} (Cross Stage Partial with 2 convolutions and feature fusion) модулиар сольсноор параметрийн тоо багасч, хэв шинжийг ялгах градиентийн урсгал сайжирсан.
    \item \textbf{Neck (Хүзүү):} PANet (Path Aggregation Network) маягийн бүтцийг ашиглан Backbone-оос ирсэн өөр өөр нягтаршлын түвшний (scale) онцлогуудтай давхаргуудыг хооронд нь холбож нэгтгэдэг. Ингэснээр жижиг бараа (жишээ нь зай) болон том хайрцагтай бараануудыг нэгэн зэрэг алдаагүй таних чадвар дээшилдэг.
    \item \textbf{Head (Толгой):} YOLOv8 нь \textit{Anchor-Free} (зангуугүй) толгойтой болсон бөгөөд хуучин шиг урьдчилаад зохиомол тэгш өнцөгтүүдийн (anchors) хэмжээг тооцдоггүйгээр зургийн аль хэсэгт объектын төв байна, түүний хил хязгаар нь хаана хүртэл байна гэдгийг (Decoupled Head - салангид толгой) байршил болон ангиллыг тус тусад нь тооцож гаргадгаар маш том өөрчлөлт орсон юм.
\end{itemize}

\textbf{Алдааны функц (Loss Function):} Уг загвар нь сургалтын үеийн алдааг тооцоолохдоо хэд хэдэн алдааны функцийн нийлбэрийг ашигладаг. Үүнд:
\begin{equation}
    Loss = \lambda_{box} L_{box} + \lambda_{cls} L_{cls} + \lambda_{dfl} L_{dfl}
\end{equation}
Энд $L_{box}$ нь CIoU (Complete Intersection over Union) loss буюу хүрээлэх хайрцгийн нарийвчлалыг хэмжинэ. $L_{cls}$ нь ангиллыг зөв таамагласан эсэхийг хэмжих Binary Cross Entropy (BCE) loss юм. Харин $L_{dfl}$ нь DFL (Distribution Focal Loss) бөгөөд объектын ирмэг хаагуур өнгөрч байгааг нарийвчлан тооцоолоход зориулагдсан.

\section{Аудит хийх алгоритм ба математик загварчлал}
Илрүүлэлт (Detection) амжилттай хийгдэж, зураг дээрээс бараануудыг олсны дараа аудитын хөдөлгүүр (Audit Engine) ажиллаж эхэлнэ. Энэхүү модулийн гол зорилго нь "Олдсон" барааг "Байх ёстой" мэдээлэлтэй автоматаар тулгаж, логик дүгнэлт гаргах явдал юм.

Бид дараах математик моделийг ашиглаж аудитыг томьёолсон болно.
Тухайн лангууд (эсвэл дэлгүүрт) бүртгэгдсэн, хүлээгдэж буй (Expected) бараануудын багцыг $E$, зургаас танигдсан (Detected) бараануудын багцыг $D$ гэж үзье.
Бараа тус бүр жижиглэн худалдаанд $item \in Items$ гэсэн ялгаатай нэртэй (жишээ нь, 'cola', 'water') танигдсан байна.

\begin{itemize}
    \item \textbf{Хүлээгдэж буй нийт тоо (Total Expected):} Дэлгүүрт нийтдээ байх ёстой барааны тоо.
    \begin{equation}
        E_{total} = \sum_{i \in Items} E_i
    \end{equation}
    Энд $E_i$ нь $i$-р барааны хүлээгдэж буй тоо хэмжээ.
    \item \textbf{Нийт таарсан зөв тоо (Total Matched):} Олдсон тоо болон хүлээгдэж байсан тооны аль багаар нь авч зөв нийлэлт үүссэн барааг тоолно. Хэрвээ байх ёстойгоос илүү их байвал, тооцоонд хүлээгдэж байсан тоогоороо л бодогдоно.
    \begin{equation}
        M_{total} = \sum_{i \in Items} \min(D_i, E_i)
    \end{equation}
    Энд $D_i$ нь зургаас илрүүлж тоолсон тухайн барааны илэрцийн тоо.
    \item \textbf{Тохирлын хувь (Match Rate):} Таарсан дүн болон хүлээгдэж байсан дүнгийн харьцаа. Энэ нь $0 \ldots 100\%$ хүртэл утга авна.
    \begin{equation}
        Rate = \left( \frac{M_{total}}{E_{total}} \right) \times 100\%
    \end{equation}
\end{itemize}

Тэгвэл аудитын эцсийн статусыг дээрх тохирлын хувьд үндэслэн автомат дүгнэлт гаргах нөхцөлт шалгагчид доорх байдлаар тодорхойлов:
\begin{description}
    \item[\textbf{PASS}:] Хэрэв $Rate = 100\%$ буюу бүх хүлээгдэж байсан бараанууд яг бодит байдал дээрхтэйгээ таарсан үед ямар нэг асуудалгүй гэж үзэж Ногоон төлөвт орно. Зөрүү буюу $Diff_i = D_i - E_i = 0$ байна.
    \item[\textbf{WARNING}:] Хэрэв $80\% \le Rate < 100\%$ буюу цөөн хэдэн бараа дутсан тохиолдолд анхааруулах Шар төлөвт орно. Лангууны хойд талд бараа нуугдсан, эсвэл түр зуур зарагдаад цөөрч урд талын мөрийг сийрэгжүүлсэн үед ийм дохио ирнэ. Менежер үүнийг хараад лангуунд бараа дүүргэлт (restocking) хийх арга хэмжээ авна.
    \item[\textbf{FAIL}:] Хэрэв $Rate < 80\%$ буюу өрөгдвөл зохих бараанууд огт тавигдаагүй, өөр төрлийн үнэд тохирохгүй бараанууд олноор орж ирсэн зэрэг ноцтой асуудалтай үед Улаан төлөв буюу бүтэлгүйтсэн гэж үзнэ. Энэ тохиолдолд төвөөс яаралтай хүн илгээж зөрчлийг арилгана.
\end{description}

\section{Сервер болон хэрэглэгчийн талын архитектурын онцлогууд (FastAPI \& Flutter)}
FastAPI нь Python хэлний төрлийн заагч (typing hints) дээр тулгуурлаж ажилладаг бөгөөд энэ нь кодчиллын үед гарах алдааг бууруулдаг. Мөн \texttt{asyncio} технологийг ашиглаж гадаад санах ой руу (MongoDB) бичих, зураг илгээж хүлээх зэрэг O/I (хүлээн авах/гаргах) үйлдлүүдийн үед серверийн ачааллыг түгжихгүйгээр бусад хэрэглэгчдийн хүсэлтийг давхар шийдвэрлэх (Concurrency) боломжтойг систем боловсруулах аргуудын үндэс болгож тавьсан.

Илүү олон ухаалаг төхөөрөмжийг нэгэн зэрэг дэмжихийн тулд бид UI бүтээхдээ Flutter-ийг сонгосон бөгөөд Flutter нь бүх элементүүдийг Виджет (Widget) гэж үздэг. Аппликейшний интерфейсийн хөгжүүлэлтэд UI ба Бизнес Логикийг хооронд нь тусгаарлах замаар (MVVM архитектур), сүлжээний алдаа, өгөгдөл шинэчлэгдэх явцуудыг хариуцдаг Provider State Management ашиглан нэг доороос удирдаж буй нь хэрэглэгчийн туршлагыг (UX) дээшлүүлэх аргачлал юм.
